%!TEX program = xelatex
% 加载配置（别动）
%!TEX program = xelatex
% 广东工业大学 本科毕业设计(论文) 标准格式
% 页边距、字体、行距、标题、目录、页码 全部符合手册
% ==============================================
\documentclass[12pt,a4paper,openany,UTF8]{ctexbook}

\usepackage{fontspec}
% 中文宋体、黑体；英文 Times New Roman
\setCJKmainfont{SimSun}
\setmainfont{Times New Roman}

% 页边距（学校要求：上30 下25 左30 右20）
\usepackage{geometry}
\geometry{
    top=30mm,
    bottom=25mm,
    left=30mm,
    right=20mm,
    includeheadfoot
}

% 1.5倍行距 + 首行缩进2字符
\usepackage{setspace}
\onehalfspacing
\setlength{\parindent}{2em}
\setlength{\parskip}{0pt}

% 标题格式（完全按手册）
\usepackage{titlesec}
% 章：三号黑体加粗
\titleformat{\chapter}[block]
    {\centering\fontsize{16pt}{18pt}\selectfont\bfseries\heiti}
    {第\thechapter 章}
    {1em}{}
% 节：小四黑体加粗
\titleformat{\section}[block]
    {\fontsize{12pt}{14pt}\selectfont\bfseries\heiti}
    {\thesection}{0.5em}{}
% 条：小四黑体
\titleformat{\subsection}[block]
    {\fontsize{12pt}{14pt}\selectfont\heiti}
    {\thesubsection}{0.5em}{}

% 页码：页脚右侧、小五号 Times New Roman
\usepackage{fancyhdr}
\pagestyle{fancy}
\fancyhf{}
\fancyfoot[R]{\fontsize{10.5pt}{12pt}\selectfont \thepage}
\renewcommand{\headrulewidth}{0pt}

% 必备宏包
\usepackage{amsmath}
\usepackage{graphicx}
\usepackage{booktabs}
\usepackage{enumitem}
\usepackage{natbib}
\bibliographystyle{unsrtnat}

% 图表按章编号（图1.1、表1.1）
\usepackage{chngcntr}
\counterwithin{figure}{chapter}
\counterwithin{table}{chapter}

\begin{document}
\begin{document}

% ===================== 封面 =====================
\thispagestyle{empty}
\begin{center}
\vspace*{3cm}

\heiti\fontsize{22pt}{24pt}\bfseries 广东工业大学

\vspace{2.5cm}

\heiti\fontsize{16pt}{18pt}\bfseries 本科毕业设计（论文）

\vspace{2cm}

\heiti\fontsize{15pt}{17pt}\bfseries
论文题目：XXXXXXXXXXXXXXXXXXXX

\vspace{2cm}

\fontsize{14pt}{16pt}\selectfont
学院：\_\_\_\_\_\_\_\_\_\_\_\_\_\_\_\_\_\_ \\
专业：\_\_\_\_\_\_\_\_\_\_\_\_\_\_\_\_\_\_ \\
学号：\_\_\_\_\_\_\_\_\_\_\_\_\_\_\_\_\_\_ \\
姓名：\_\_\_\_\_\_\_\_\_\_\_\_\_\_\_\_\_\_ \\
指导教师：\_\_\_\_\_\_\_\_\_\_\_\_\_\_\_\_\_\_

\vspace{3cm}

2026 年 5 月
\end{center}
\clearpage

% ===================== 中文摘要 =====================
\pagenumbering{Roman}
\setcounter{page}{1}

\chapter*{\heiti 摘\quad 要}
\addcontentsline{toc}{chapter}{摘要}

（这里写中文摘要，约400字）

\vspace{1em}
\noindent\textbf{关键词：}XXX；XXX；XXX；XXX

\clearpage

% ===================== 英文摘要 =====================
\chapter*{\bfseries ABSTRACT}
\addcontentsline{toc}{chapter}{ABSTRACT}

（英文摘要）

\vspace{1em}
\noindent\textbf{Keywords：}XXX；XXX；XXX

\clearpage

% ===================== 目录 =====================
\chapter*{\heiti 目\quad 录}
\tableofcontents
\clearpage

% ===================== 正文页码（阿拉伯数字） =====================
\pagenumbering{arabic}
\setcounter{page}{1}

% ===================== 第1章 绪论 =====================
\chapter{绪论}
\section{研究背景与意义}
\subsection{研究背景}
这里写正文内容，自动小四宋体、1.5倍行距、首行缩进。

\subsection{研究意义}
本文主要完成以下工作：
\begin{enumerate}[label=(\arabic*)]
    \item 研究XXX
    \item 设计XXX
    \item 实现XXX
\end{enumerate}

\section{国内外研究现状}
国内外学者在XXX领域开展了大量研究\cite{example1}。

% ===================== 第2章 相关理论 =====================
\chapter{相关理论与技术}
\section{Spark Streaming 原理}
DStream是离散流，本质是RDD序列。

% ===================== 公式示例 =====================
\section{公式写法}
单词统计计算公式如下：
\begin{equation}\label{eq:count}
    Count(w) = \sum_{i=1}^N 1
\end{equation}
由式(\ref{eq:count})可知XXX。

% ===================== 图片示例 =====================
\section{图片写法}
\begin{figure}[htbp]
    \centering
    \includegraphics[width=0.8\textwidth]{figures/arch.png}
    \caption{Spark Streaming 架构图}
    \label{fig:arch}
\end{figure}
图\ref{fig:arch}为系统整体架构。

% ===================== 表格示例（三线表） =====================
\section{表格写法}
\begin{table}[htbp]
    \centering
    \caption{实验环境}
    \label{tab:env}
    \begin{tabular}{cc}
        \toprule
        项目 & 配置 \\
        \midrule
        系统 & CentOS 7 \\
        Spark版本 & 2.4.0 \\
        \bottomrule
    \end{tabular}
\end{table}

% ===================== 结论 =====================
\chapter{结论与展望}
本文完成XXX，未来可进一步研究XXX。

% ===================== 参考文献 =====================
\chapter*{\heiti 参考文献}
\addcontentsline{toc}{chapter}{参考文献}
\bibliography{ref}

% ===================== 致谢 =====================
\chapter*{\heiti 致\quad 谢}
\addcontentsline{toc}{chapter}{致谢}
感谢指导教师XXX的悉心指导，感谢实验室同学的帮助。

% ===================== 附录 =====================
\chapter*{\heiti 附\quad 录}
\addcontentsline{toc}{chapter}{附录}
附录A：核心代码

\end{document}